\subsection{Related work}
\label{sec:related-work}

%% We briefly overview related work, acknowledging that 
To make our theoretical model tractable, we capture only a few of the complexities inherent in
dataset construction. Label aggregation strategies often attempt to
evaluate labeler uncertainty, dating to~\citet{DawidSk79}, who study
labeler uncertainty estimation to overcome noise in
clinical patient measurements. With the rise of
crowdsourcing, such reliability models have
attracted interest, with approaches for optimal budget
allocation~\cite{KargerOhSh14}, addressing untrustworthy or malicious
labelers~\cite{CharikarStVa17}, and more broadly a line of applied~\cite{DengDoSoLiLiFe09, RussakovskyDeSuKrSaMaHuKaKhBeBeFe15} and theoretical work
studying crowd labeling and aggregation~\cite{WhitehillWuBeMoRu09,
  WelinderBrPeBe10, RatnerBaEhFrWuRe17}.
%%, with substantial
%% applications~\cite{DengDoSoLiLiFe09, RussakovskyDeSuKrSaMaHuKaKhBeBeFe15}.

Many of these focus, however, on obtaining a single
trustworthy label for each example. Our work takes a
different perspective.  First, we target statistical analysis for
the full learning pipeline---understanding the landscape of the
learning problem with multiple labels for each example---as opposed to
obtaining only clean labels. Moreover, we argue for
calibrating the resulting predictive model, which cleaned
labels necessarily impede. We therefore adopt a perspective similar to
the applied works~\citep{PetersonBaGrRu19,PlataniosAlXiMi20},
which highlight how incorporating human
uncertainty into learning pipelines can make classification ``more
robust''~\citep{PetersonBaGrRu19}.  Instead of splitting the learning
procedure into two phases, where the first aggregates labels and the second
trains, we use non-aggregated labels throughout
learning. \citet{PlataniosAlXiMi20} directly
model labeler reliability via a richer model than our stylized
scenarios, but our
approach allows theoretically precise description of the
methods' limiting
behavior.


Our results complement a growing body of work in machine learning.
%
The Neural Information Processing Systems conference
includes a Datasets and Benchmarks track, as they
``are crucial for the development of machine learning
methods''~\cite{BeygelzimerLiVaDa21},
%
%% As a particular examplar, \citet{GadreEtAl23} introduce DataComp,
%% a testbed for experimental work on datasets,
building off the long history
of dataset-driven development in machine learning~\cite{HardtRe22}.
%
We attempt to lay statistical foundations for this line of work.

%We mention in passing that our consistency results rely on
%distributional assumptions on the covariates $X$, for example,
%Gaussianity. That such technical conditions appear may be familiar from work,
%for example, in single-index or nonlinear measurement models~\cite{PlanVe16,
%  PlanVe13}. In such settings, one assumes $\E[Y \mid X] = f(\<\theta\opt,
%X\>)$ for an unknown increasing $f$, and it is essential that $\E[Y X]
%\propto \theta\opt$ to allow estimation of $\theta\opt$; we leverage similar
%techniques.

\paragraph{Notation}
For a matrix $M$, $\norm{M}$ denotes its spectral norm and $M^\dagger$ its
Moore-Penrose pseudo inverse.
%
For a unit vector $u$, $\proj_{u}^\perp \defeq
I - u u^\top$ projects to the space orthogonal to $\mathrm{span}\{u\}$.
%
We use the empirical process notation $\Prb Z = \int z d\Prb(z)$ and $\cd$
denotes convergence in distribution.
%
We write $f(x) \asymp g(x)$ for $x \in \R_+$ if there exist
numerical constants $0 < c_1, c_2 < \infty$ and $x_0 \geq 0$ such that $c_1
|g(x)| \leq |f(x)| \leq c_2|g(x)|$ for $x \geq x_0$,
and
$c = o_m(1)$ means $c \to 0$ as $m \to \infty$.

